% Quatrième de couverture provisoire : remplacée automatiquement par
% lastpage.pdf si ce fichier (modèle officiel ESPRIT) est ajouté au dossier.
\clearpage
\thispagestyle{empty}
\noindent\textbf{\large Résumé}\\[0.3cm]
Ce rapport présente le travail réalisé au sein d'Antigone dans le cadre de notre projet de fin d'études. Il porte sur la conception et le développement de Sellio, une plateforme transactionnelle intelligente de type SaaS qui permet à des commerçants de cosmétiques et de vêtements de créer et d'exploiter leur boutique en ligne. La plateforme repose sur une architecture microservices (Spring Boot, Spring Cloud Gateway, PostgreSQL) et des interfaces React. Elle intègre la vérification d'identité des commerçants, trois gabarits de vitrine personnalisables, le paiement par Stripe, des outils marketing et de fidélisation, un essayage virtuel en temps réel fondé sur l'API Decart, ainsi qu'un suivi comportemental exploité par apprentissage automatique : recommandation hybride, segmentation des visiteurs et prédiction de l'attrition, évaluées sur des jeux de données publics réels. Le projet a été conduit selon la méthode Scrum.

\vspace{0.3cm}
\noindent\textbf{Mots clés :} e-commerce, SaaS, microservices, Spring Boot, React, Stripe, essayage virtuel, apprentissage automatique, système de recommandation, attrition, Scrum.

\vspace{1.2cm}
\noindent\textbf{\large Abstract}\\[0.3cm]
This report presents the work carried out at Antigone as our end-of-studies project. It covers the design and development of Sellio, an intelligent transactional SaaS platform that lets cosmetics and clothing merchants create and run their own online store. The platform is built on a microservices architecture (Spring Boot, Spring Cloud Gateway, PostgreSQL) with React front ends. It includes merchant identity verification, three customisable storefront templates, Stripe payments, marketing and loyalty tools, real-time virtual try-on based on the Decart API, and behavioural tracking used by machine-learning models: hybrid recommendation, visitor segmentation and churn prediction, evaluated on real public datasets. The project followed the Scrum framework.

\vspace{0.3cm}
\noindent\textbf{Keywords:} e-commerce, SaaS, microservices, Spring Boot, React, Stripe, virtual try-on, machine learning, recommender system, churn, Scrum.
